
\section{\sotopia interaction environment}
\label{sec:sotopia_characters_relationships}
To address the challenge of evaluating social intelligence interactively, we seek an environment with the following desiderata:  (1) \textit{Realistic}: this is to evaluate and understand artificial agents' behavior under realistic scenarios; (2) \textit{Mixed utilities}: human motives are often driven by both explicit and implicit incentives, and the environment should be able to evaluate the agents' performance on multiple dimensions; (3) \textit{Open-ended}: to support large-scale simulation and evaluation, the environment should be able to produce new tasks satisfying the previous two desiderata procedurally, without heavy human intervention.

In this section, we introduce \sotopia and explain why \sotopia is well-suited for interactive evaluation of social intelligence. The task space includes realistic scenarios, characters, and relationships which are automatically generated with manual inspection (\S\ref{sec:task_space}). An episode includes the interaction between agents role-playing different characters who each perform actions (e.g. \texttt{speak("Hello Bob!")}, \texttt{smile and nod}, and \texttt{call 911}) to achieve social goals drawn from the task space (\S\ref{sec:episodes}). We direct readers to Appendix \ref{app:formal_formulation} for a formal definition of the \sotopia environment. 


\subsection{Task space}
\label{sec:task_space}
In this paper, we consider tasks that involve two agents, but \sotopia is more general and could support the interaction among more than two agents. 
A task in \sotopia is the combination of a \emph{scenario context}, \emph{characters}, and their \emph{social goals}, providing the background of the interaction. Each episode consists of multiple turns of interaction between agents. In this paper, we focus on locally-consistent social goals within a relatively short timespan in single episodes, despite that in the real world, people's social goals are consistently changing from time to time. Note that agents have different observations for the same task: each agent can observe the scenario, their own social goal, and their own character profile. Other agents' social goals are invisible and other agents' character profiles are partially observable, depending on the relationship between the agents. 

\paragraph{Complexity of task space}
The combinations of a scenario context, social goals, characters, and their relationships can shape the space of the optimal behaviors of agents. Consider a persuasion task, ``asking the romantic partner to stop texting during FaceTime.'' If a romantic partner values conformity, one good way for an agent to reach this goal is to discuss the problem from a social norm perspective; however, if a romantic partner is particularly caring and good at understanding feelings, it might be better to express subjective emotion. \emph{Interaction partner's policy} also heavily influences the optimal behaviors. Consider another task illustrated in Figure \ref{fig:teaser}, ``\textsl{selling BMW Z3 for no less than \$3,400}''. If the buyer gives a high offer, the seller might want to exploit the buyer's eagerness to buy the car and ask for a higher price; while if the buyer gives a low-ball offer, the seller could give reasons why the car is worth more than that or threaten to walk away. When more information (e.g.\ about personality, decision-making styles, or occupation) is known before the interaction, the seller and buyer could use that knowledge to adjust their strategies as well. 
The cross-product of the diverse spaces of scenario context, social goals, characters, relationship profiles, and other players' policies creates a large task space that poses not only a realistic challenge but also an opportunity to evaluate and develop social intelligence in artificial agents. For the rest of this subsection, we will present the design and generation of each axis of the task space. 

\paragraph{Characters}
As mentioned above, the design of character profiles should include several attributes that would influence decision-making. We consider the following ones (inspired by \citet{wang-etal-2019-persuasion}): name, gender, age, occupation, pronouns, personality traits \citep{goldberg1992-jf}, moral values \citep{graham2011-aw}, Schwartz personal values \citep{cieciuch_davidov_2012}, and decision-making style \citep{hamilton2016}, which are generated through leveraging GPT-4 \citep{openai2023gpt4}. To give the conversations more background, after generating the above attributes, we prompt GPT-4 to generate secret and public information.
Two examples of characters are shown in Figure \ref{fig:teaser}.
It should be noted that, although we generated a diverse set of characters, this is still a small portion of the possible character space. Our analysis focuses on 40 characters generated in the aforementioned fashion, and future research using \sotopia can easily generate an expanded character set.

\paragraph{Relationships}
Relationships in \sotopia have the following effects: (1) scenarios often have \textit{relationship constraints}; for example, a family relationship is required for a family dinner scenario, but not for a scenario involving finding mutual friends at a party; (2) different relationships influence an agent's observation of the profiles of other agents during interactions; for example, a stranger may not have knowledge about another agent's occupation, while a romantic partner may know the other agent's personality. To make sampling characters easier for (1) and controlling the interaction context easier for (2), we consider five types of relationships: \textit{family}, \textit{friend}, \textit{romantic}, \textit{acquaintance}, and \textit{stranger}. Refer to Appendix \ref{sec:limitations} for the limitations of this approach and potential extensions.

We will discuss how (1) is performed in the following paragraphs, while for (2), we created a rule-based mechanism to determine whether the parts of the profiles are visible to the other agent. If two agents are in family, friends, or romantic relationships, they can see everything on each other's profile except for secrets. Two acquaintances can see the name, occupation, gender pronouns, and public info on each other's profile. Two strangers can see nothing on each other's profile. Similar to characters, we prompt GPT-4 \citep{openai2023gpt4} to automatically generate relationships based on the character pool and manually validate relationships for consistency. 

\paragraph{Scenarios}
We consider scenarios where the agents have both shared and private information about the social task.
The shared information is the scenario context: the location, time and other shared information of the social interaction, e.g.\ ``\textsl{One person is selling an antique chair for \$100 on his patio. Another person is interested in this chair.}'' The private information is the social goals which are only visible to the respective agents, e.g.\ ``\textsl{Your goal is to buy the chair for \$80.}'' is only visible to the buyer agent, while ``\textsl{Your goal is to sell the chair for \$90.}'' is only visible to the seller agent. However, the as mentioned above combination of scenarios and characters is not arbitrary, since scenarios often imply constraints for the agents. We call this kind of constraint \emph{scenario constraints}. In this paper, we mainly consider \emph{relationship constraints} which determines the types of relationships between the sampled characters. Similar to characters and relationships, scenarios, including context, goals, and constraints are generated through prompting GPT-4 \citep{openai2023gpt4}. To generate high-quality scenarios with enough coverage of different types of social interactions (as shown in Figure \ref{fig:teaser}), we randomly sample data from previous datasets, including \citealt{forbes-etal-2020-social, sap-etal-2019-social, lewis-etal-2017-deal,ziems-etal-2023-normbank, he-etal-2018-decoupling, he-etal-2017-learning}, and use them in the prompts to ``inspire'' GPT-4. The authors manually validate and make necessary changes to all of the generated scenarios and remove 10\% of scenarios according to \ref{app_sec:annotation_guidelines}. 

\subsection{\sotopia episodes}
\label{sec:episodes}
During the interaction, models and humans are given the social context, a character profile and a corresponding social goal. 
We will call these models and humans with characters and goals \emph{agents}, which take turns (in a round-robin fashion, i.e. Agent 1 acts first and then Agent 2 acts and so on) to perform actions in an \emph{episode}. 
At their own turn, the agent can choose to \texttt{speak}, use \texttt{non-verbal communication} (e.g., hug or smile in Figure \ref{fig:qualitative_example_hug_demo}), or take a \texttt{physical action} (e.g., 
play music in Figure \ref{fig:qualitative_example_play_music_demo}), which are all important components of social interactions \citep{De_Stefani2019-ih}. 
Once an agent chooses one of these three discrete action categories, the agent then generates a specific action, i.e. what to say, what gesture to make, etc., in text form.
Outside of the three actions, the agent can also choose to do nothing (\texttt{none}) to express silence or allow another agent to finish, or choose to \texttt{leave} to end the episode. We set the limit of the turns to 20, as we found humans normally can finish most of the tasks in 20 turns. An episode ends either because one of the agents chooses to leave, or it reaches the limit of turns. An example episode is shown in Figure \ref{fig:teaser}.